\textbf{Это образец заполнения.} Автор, группа и дата ниже вымышлены.
Свою лабораторную работу оформляйте по наблюдениям на собственном
компьютере; отправляйте преподавателю отчёт в \textbf{PDF или HTML}, а
исходные файлы сохраняйте отдельно для защиты.

\textbf{Автор:} Иванов Иван.\\
\textbf{Группа:} учебная группа ПИ-000.\\
\textbf{Дата:} 1 октября 2026 года.\\
\textbf{Тема:} проверка суммы двух чисел и исправление логической
ошибки.

\hypertarget{ux446ux435ux43bux44c-ux438-ux443ux441ux43bux43eux432ux438ux44f}{%
\section{Цель и
условия}\label{ux446ux435ux43bux44c-ux438-ux443ux441ux43bux43eux432ux438ux44f}}

Создать программу на C, вычисляющую сумму двух целых чисел. Перед
запуском предсказать результат, сравнить его с выводом, исследовать
причину расхождения в GDB и проверить исправление. Исходные данные:
\texttt{a\ =\ 7}, \texttt{b\ =\ 5}. Ожидаемый результат:
\textbf{\texttt{7\ +\ 5\ =\ 12}}.

\hypertarget{ux441ux440ux435ux434ux430-ux438-ux438ux441ux445ux43eux434ux43dux44bux439-ux444ux430ux439ux43b}{%
\section{Среда и исходный
файл}\label{ux441ux440ux435ux434ux430-ux438-ux438ux441ux445ux43eux434ux43dux44bux439-ux444ux430ux439ux43b}}

Для воспроизведения использовалась среда Linux (Debian 12), GCC 12.2 и
GDB 13.1 с параметрами \texttt{-g\ -O0}. Я создал каталог
\texttt{\textasciitilde{}/aiya/demo-sum} и набрал \texttt{sum.c} в
редакторе nano. Исходная ключевая строка была:

\begin{verbatim}
int result = a - b;
\end{verbatim}

Остальной код выводил результат через
\texttt{printf("Сумма:\ \%d\textbackslash{}n",\ result);} и возвращал
\texttt{0} из \texttt{main}. Полный исходник до исправления показан в
демонстрационной лабораторной; файл программы я сохранил отдельно от
отчёта.

\hypertarget{ux43fux435ux440ux432ux44bux439-ux437ux430ux43fux443ux441ux43a-ux438-ux440ux430ux441ux445ux43eux436ux434ux435ux43dux438ux435}{%
\section{Первый запуск и
расхождение}\label{ux43fux435ux440ux432ux44bux439-ux437ux430ux43fux443ux441ux43a-ux438-ux440ux430ux441ux445ux43eux436ux434ux435ux43dux438ux435}}

Из рабочего каталога я выполнил:

\begin{verbatim}
gcc -std=c11 -Wall -Wextra -g -O0 sum.c -o sum
./sum
echo $?
\end{verbatim}

GCC завершил сборку без предупреждений. Программа напечатала
\textbf{\texttt{Сумма:\ 2}}; следующий \texttt{echo\ \$?} показал
\texttt{0}. До запуска я ожидал \texttt{Сумма:\ 12}. Следовательно,
программа завершилась успешно с точки зрения её собственного кода
возврата, но вычислила \textbf{не ту операцию}. Отсутствие
предупреждений компилятора не означает, что логика задачи верна.

\hypertarget{ux43eux442ux43bux430ux434ux43aux430}{%
\section{Отладка}\label{ux43eux442ux43bux430ux434ux43aux430}}

Я запустил \texttt{gdb\ -q\ ./sum} и ввёл команды в приглашении
\texttt{(gdb)}:

\begin{verbatim}
break sum.c:7
run
print a
print b
next
print result
quit
\end{verbatim}

Остановка была \textbf{перед} вычислением \texttt{result}. Наблюдения:
\texttt{print\ a} → \texttt{7}, \texttt{print\ b} → \texttt{5}; после
\texttt{next} команда \texttt{print\ result} → \texttt{2}. Числа в
обозначениях ответов GDB (\texttt{\$1}, \texttt{\$2} и т. п.) относятся
к истории отладчика и не являются адресами переменных. В исходнике
стояло вычитание вместо требуемого сложения; проблема не была связана с
вводом данных или сборкой.

\hypertarget{ux438ux441ux43fux440ux430ux432ux43bux435ux43dux438ux435-ux438-ux43fux43eux432ux442ux43eux440ux43dux430ux44f-ux43fux440ux43eux432ux435ux440ux43aux430}{%
\section{Исправление и повторная
проверка}\label{ux438ux441ux43fux440ux430ux432ux43bux435ux43dux438ux435-ux438-ux43fux43eux432ux442ux43eux440ux43dux430ux44f-ux43fux440ux43eux432ux435ux440ux43aux430}}

Исправил ровно одну строку и \textbf{пересобрал} программу:

\begin{verbatim}
// было: int result = a - b;
int result = a + b;
\end{verbatim}

Для каждой дополнительной пары менял \texttt{a} и \texttt{b} в
\texttt{sum.c}, снова выполнял указанную выше команду GCC и запускал
\texttt{./sum}. Получил:

\begin{longtable}[]{@{}
  >{\raggedleft\arraybackslash}p{(\columnwidth - 8\tabcolsep) * \real{0.0700}}
  >{\raggedleft\arraybackslash}p{(\columnwidth - 8\tabcolsep) * \real{0.0700}}
  >{\raggedright\arraybackslash}p{(\columnwidth - 8\tabcolsep) * \real{0.2600}}
  >{\raggedright\arraybackslash}p{(\columnwidth - 8\tabcolsep) * \real{0.3400}}
  >{\raggedleft\arraybackslash}p{(\columnwidth - 8\tabcolsep) * \real{0.2200}}@{}}
\toprule
\begin{minipage}[b]{\linewidth}\raggedleft
\texttt{a}
\end{minipage} & \begin{minipage}[b]{\linewidth}\raggedleft
\texttt{b}
\end{minipage} & \begin{minipage}[b]{\linewidth}\raggedright
Ожидалось
\end{minipage} & \begin{minipage}[b]{\linewidth}\raggedright
Наблюдалось
\end{minipage} & \begin{minipage}[b]{\linewidth}\raggedleft
Код завершения
\end{minipage} \\
\midrule
\endhead
7 & 5 & \texttt{Сумма:\ 12} & \texttt{Сумма:\ 12} & 0 \\
0 & 5 & \texttt{Сумма:\ 5} & \texttt{Сумма:\ 5} & 0 \\
-3 & 5 & \texttt{Сумма:\ 2} & \texttt{Сумма:\ 2} & 0 \\
\bottomrule
\end{longtable}

В конце вернул значения \texttt{a\ =\ 7}, \texttt{b\ =\ 5} и снова
собрал программу. Файл \texttt{sum.c} и исполняемый файл \texttt{sum}
сохранил для повторного показа.

\hypertarget{ux432ux44bux432ux43eux434}{%
\section{Вывод}\label{ux432ux44bux432ux43eux434}}

Сравнение с заранее записанным ожиданием помогло обнаружить ошибку,
которую GCC не распознал: программа вычитала \texttt{b} из \texttt{a}
вместо сложения. Просмотр переменных и выполнение строки в GDB
подтвердили причину результата \texttt{2}. После исправления и
пересборки все три контрольные проверки совпали с ожидаемыми значениями.
При подготовке собственного отчёта важно записывать \textbf{и неудачную
попытку}, и подтверждение исправления --- с командами и наблюдениями.
